\section{System Requirements (Requisiti di Sistema)}

I requisiti di sistema formalizzano in dettaglio le capacità funzionali, i vincoli qualitativi e le regole di business che la piattaforma software \textbf{MyAma} deve soddisfare, in conformità alle linee guida dello standard IEEE 830-1998.

% =========================================================================
% 4.1 REQUISITI FUNZIONALI
% =========================================================================
\subsection{Requisiti Funzionali}
I requisiti funzionali descrivono i servizi e le azioni che il sistema deve eseguire in risposta a specifici stimoli o interazioni.

\begin{center}
\renewcommand{\arraystretch}{1.3}
\begin{longtable}{|p{2cm}|p{8.2cm}|p{4.8cm}|}
\hline
\rowcolor{tableheader} \textbf{ID} & \textbf{Descrizione del Requisito} & \textbf{Use Case di Riferimento} \\ \hline
\endfirsthead

\hline
\rowcolor{tableheader} \textbf{ID} & \textbf{Descrizione del Requisito} & \textbf{Use Case di Riferimento} \\ \hline
\endhead

\textbf{RF-01} & Il sistema deve consentire la libera consultazione delle categorie di rifiuti ingombranti e delle relative tariffe di smaltimento. & UC-U01 / Visitatore \\ \hline

\textbf{RF-02} & Il sistema deve permettere la ricerca delle sedi AMA e isole ecologiche attive tramite inserimento del CAP. & UC-U01 / Visitatore \\ \hline

\textbf{RF-03} & Il sistema deve consentire la registrazione di un nuovo profilo Cittadino previa validazione dei dati anagrafici e recapiti. & UC-U02 / Visitatore \\ \hline

\textbf{RF-04} & Il sistema deve consentire al Cittadino di richiedere un ritiro a domicilio indicando tipologia, descrizione e quantità dei beni. & UC-C01 / Cittadino \\ \hline

\textbf{RF-05} & Il sistema deve validare la copertura territoriale verificando che il CAP indicato appartenga alle zone servite da AMA. & UC-C01 / Sistema \\ \hline

\textbf{RF-06} & Il sistema deve calcolare gli slot temporali disponibili tenendo conto della portata residua dei veicoli e dei turni autisti. & UC-C01 / Sistema \\ \hline

\textbf{RF-07} & Il sistema deve consentire il caricamento opzionale di fotografie del rifiuto per agevolare la stima volumetrica. & UC-C01 / Cittadino \\ \hline

\textbf{RF-08} & Il sistema deve permettere la prenotazione di un conferimento diretto presso un centro di raccolta autorizzato. & UC-C02 / Cittadino \\ \hline

\textbf{RF-09} & Il sistema deve filtrare e mostrare le sole sedi abilitate alla ricezione della specifica tipologia di rifiuto indicata. & UC-C02 / Sistema \\ \hline

\textbf{RF-10} & Il sistema deve generare un Pass di Conferimento con codice identificativo univoco e QR code all'atto della conferma. & UC-C02 / Cittadino \\ \hline

\textbf{RF-11} & Il sistema deve consentire al Cittadino di visualizzare lo stato di avanzamento e lo storico di tutte le proprie prenotazioni. & UC-C03 / Cittadino \\ \hline

\textbf{RF-12} & Il sistema deve consentire la cancellazione autonoma della prenotazione se richiesta con almeno 24 ore di anticipo. & UC-C03 / Cittadino \\ \hline

\textbf{RF-13} & Il sistema deve consentire al Cittadino di rilasciare una valutazione con punteggio e commento sui servizi completati. & UC-C04 / Cittadino \\ \hline

\textbf{RF-14} & Il sistema deve fornire all'Autista l'elenco ordinato e l'itinerario dei ritiri a domicilio assegnati per il proprio turno. & UC-A01 / Autista \\ \hline

\textbf{RF-15} & Il sistema deve mostrare all'Autista i dettagli del carico e la percentuale di capienza progressivamente occupata sul mezzo. & UC-A01 / Autista \\ \hline

\textbf{RF-16} & Il sistema deve consentire all'Autista di registrare l'esito del ritiro (\textit{Completato, Assente, Non conforme}) con note operative. & UC-A02 / Autista \\ \hline

\textbf{RF-17} & Il sistema deve aggiornare in tempo reale lo stato della prenotazione dopo la registrazione dell'esito da parte dell'autista. & UC-A02 / Sistema \\ \hline

\textbf{RF-18} & Il sistema deve consentire all'Autista di avviare il contatto o inviare notifiche di arrivo al cittadino. & UC-A03 / Autista \\ \hline

\textbf{RF-19} & Il sistema deve mostrare all'Operatore di Sede la pianificazione giornaliera degli accessi al centro di raccolta. & UC-O01 / Operatore \\ \hline

\textbf{RF-20} & Il sistema deve consentire all'Operatore di validare il Pass di Conferimento tramite scansione ottica o codice manuale. & UC-O02 / Operatore \\ \hline

\textbf{RF-21} & Il sistema deve consentire all'Operatore di registrare l'avvenuto scarico del rifiuto conforme o la non conformità riscontrata. & UC-O02 / Operatore \\ \hline

\textbf{RF-22} & Il sistema deve consentire all'Amministratore di Sede di gestire l'anagrafica del personale assegnato alla sede. & UC-S01 / Amm. Sede \\ \hline

\textbf{RF-23} & Il sistema deve permettere la pianificazione dei turni settimanali per autisti e operatori di sede. & UC-S01 / Amm. Sede \\ \hline

\textbf{RF-24} & Il sistema deve consentire la gestione della flotta veicoli della sede con impostazione di portata e stato di disponibilità. & UC-S02 / Amm. Sede \\ \hline

\textbf{RF-25} & Il sistema deve consentire la configurazione degli orari del centro di raccolta e l'associazione dei CAP serviti. & UC-S03 / Amm. Sede \\ \hline

\textbf{RF-26} & Il sistema deve consentire all'Amministratore Generale di creare, abilitare e revocare gli account degli amministratori di sede. & UC-G01 / Amm. Generale \\ \hline

\textbf{RF-27} & Il sistema deve generare report e cruscotti statistici aggregate sui volumi di rifiuti gestiti e tassi di successo dei servizi. & UC-G02 / Amm. Generale \\ \hline
\end{longtable}
\end{center}

% =========================================================================
% 4.2 REQUISITI NON FUNZIONALI
% =========================================================================
\newpage
\subsection{Requisiti Non Funzionali}
I requisiti non funzionali definiscono i vincoli qualitativi di prestazione, affidabilità, usabilità, sicurezza e compatibilità del software.

\begin{center}
\renewcommand{\arraystretch}{1.3}
\begin{longtable}{|p{2.2cm}|p{3cm}|p{10cm}|}
\hline
\rowcolor{tableheader} \textbf{ID} & \textbf{Categoria} & \textbf{Descrizione del Requisito e Metrica di Misura} \\ \hline
\endfirsthead

\hline
\rowcolor{tableheader} \textbf{ID} & \textbf{Categoria} & \textbf{Descrizione del Requisito e Metrica di Misura} \\ \hline
\endhead

\textbf{RNF-01} & Prestazioni & Il tempo di risposta per il calcolo delle disponibilità e verifica CAP non deve superare i 2.0 secondi al 95° percentile sotto carico di 200 richieste concorrenti. \\ \hline

\textbf{RNF-02} & Disponibilità & Il sistema deve garantire una disponibilità del servizio del 99.5\% su base mensile nell'intervallo orario 07:00 -- 22:00. \\ \hline

\textbf{RNF-03} & Usabilità & L'interfaccia di prenotazione per il cittadino deve essere completabile in non più di 4 passaggi sequenziali intuitivi, con supporto completo per dispositivi mobili (\textit{responsive}). \\ \hline

\textbf{RNF-04} & Sicurezza & I dati in transito devono essere protetti mediante crittografia TLS 1.3; le password devono essere archiviate con algoritmo di hashing sicuro (bcrypt con work factor $\ge 12$). \\ \hline

\textbf{RNF-05} & Integrità Dati & Le operazioni di prenotazione e allocazione slot devono essere eseguite sotto transazioni ACID per prevenire sovrapposizioni e overbooking sui mezzi. \\ \hline

\textbf{RNF-06} & Portabilità & L'applicazione web deve risultare pienamente compatibile con i principali browser moderni (Google Chrome $\ge 110$, Mozilla Firefox $\ge 110$, Safari $\ge 16$, Microsoft Edge $\ge 110$). \\ \hline
\end{longtable}
\end{center}

% =========================================================================
% 4.3 REQUISITI DI DOMINIO
% =========================================================================
\subsection{Requisiti di Dominio}
I requisiti di dominio formalizzano le regole operative e normative dell'azienda municipalizzata AMA:

\begin{itemize}[leftmargin=*]
    \item \textbf{RD-01 (Vincolo di Portata Massima Veicolo):} Per ciascun veicolo impiegato nei ritiri a domicilio, la sommatoria dei pesi stimati dei rifiuti prenotati non può in nessun caso superare la portata massima omologata del mezzo ($\sum_{i} P_i \le \text{PortataMax}$).
    
    \item \textbf{RD-02 (Competenza Territoriale del Bacino CAP):} Un autista e un veicolo possono essere assegnati esclusivamente a ritiri ricadenti all'interno dei CAP formalmente associati al distretto operativo di appartenenza.
    
    \item \textbf{RD-03 (Normativa Smaltimento RAEE e Rifiuti Speciali):} Il conferimento di apparecchiature elettriche ed elettroniche (RAEE) e rifiuti speciali è consentito unicamente presso le sedi AMA provviste di specifiche autorizzazioni ambientali e contenitori dedicati.
    
    \item \textbf{RD-04 (Finestra Temporale di Cancellazione):} Il cittadino può procedere alla cancellazione autonoma della prenotazione con liberazione delle risorse fino a un massimo di 24 ore prima dell'inizio dello slot orario prenotato.
\end{itemize}

% =========================================================================
% 4.4 VERIFICABILITÀ DEI REQUISITI
% =========================================================================
\subsection{Verificabilità dei Requisiti}
In accordo con i criteri di ingegneria del software, ogni requisito deve essere controllabile e dimostrabile tramite specifiche procedure di collaudo.

\begin{center}
\renewcommand{\arraystretch}{1.3}
\begin{longtable}{|p{2cm}|p{5.5cm}|p{7.8cm}|}
\hline
\rowcolor{tableheader} \textbf{ID Req.} & \textbf{Metodo di Verifica} & \textbf{Criterio di Accettazione e Procedura di Test} \\ \hline
\endfirsthead

\hline
\rowcolor{tableheader} \textbf{ID Req.} & \textbf{Metodo di Verifica} & \textbf{Criterio di Accettazione e Procedura di Test} \\ \hline
\endhead

\textbf{RF-05} & Test Funzionale & Inserimento di CAP appartenenti ed esterni alla copertura AMA; il sistema blocca i CAP non coperti e fornisce opportuno messaggio di redirect al conferimento in sede. \\ \hline

\textbf{RF-06} & Test di Integrazione & Simulazione di prenotazioni consecutive fino a saturazione portata veicolo; il sistema inibisce la selezione dello slot quando il carico residuo è inferiore al peso del rifiuto inserito. \\ \hline

\textbf{RF-12} & Test di Limite Temporale & Tentativo di cancellazione a $T - 25\text{h}$ (operazione consentita con successo) e a $T - 23\text{h}$ (operazione bloccata con notifica di vincolo temporale). \\ \hline

\textbf{RNF-01} & Stress / Load Test & Generazione di 200 richieste concorrenti di ricerca disponibilità tramite Apache JMeter; tempo medio di risposta registrato $\le 2.0\text{s}$. \\ \hline

\textbf{RNF-04} & Security Audit & Scansione automatizzata delle vulnerabilità web (OWASP ZAP) e verifica conformità protocollo TLS 1.3 e hashing password. \\ \hline

\textbf{RD-01} & Test sui Vincoli Logici & Inserimento di rifiuti con peso complessivo superiore alla portata del mezzo assegnato; verifica del blocco transazione da parte del controllore logico. \\ \hline
\end{longtable}
\end{center}
